%% Copyright 2026 Petr Socha
%%
%% This work may be distributed and/or modified under the
%% conditions of the LaTeX Project Public License, either version 1.3c
%% of this license or (at your option) any later version.
%% The latest version of this license is in
%%   https://www.latex-project.org/lppl.txt
%% and version 1.3c is part of all distributions of LaTeX
%% version 2008 or later.
%%
%% This work has the LPPL maintenance status `maintained'.
%% The Current Maintainer of this work is Petr Socha.
%% Maintenance and bug reports: https://github.com/petrsocha/pslib-latex/issues
%%
%% The components of this work are listed in MANIFEST.txt.
%% The school logo is a separately licensed third-party asset.

% ============================================================================
% TESTOVÁNÍ -- VZORCE, MĚRNÉ JEDNOTKY A VÝSLEDKY
% ============================================================================
\chapter{Testování}
\label{chap:testovani}

\section{Matematický zápis}

V~odborné práci lze uvádět i~matematické vztahy. Například Ohmův zákon
vyjadřuje závislost proudu $I$ na napětí $U$ a~odporu $R$:

% Samostatný vzorec v prostředí equation se čísluje automaticky.
% Odkaz na vzorec funguje stejně jako odkaz na obrázek nebo tabulku.
\begin{equation}
  I = \frac{U}{R}
  \label{eq:ohm}
\end{equation}

Pro jednoduchý obvod podle \autoref{eq:ohm} s~$U=5\,\mathrm{V}$
a~$R=1\,\mathrm{k}\Omega$ vychází $I=5\,\mathrm{mA}$.

\section{Odborná přednáška jako zdroj}

% Skripta a přednáškové materiály z VŠ lze citovat stejně jako jiné zdroje.
% V literatura.bib je ukázka veřejně dostupné přednášky z MIT.
Vysokoškolské přednáškové poznámky mohou sloužit jako zdroj například
při vysvětlování filtrů. Materiál MIT \emph{Lecture 18: Filters}
popisuje mimo jiné dolní a~horní propust~\cite{mitfilters}.

U~jednoduché RC dolní propusti lze mezní frekvenci vypočítat ze vztahu
\begin{equation}
  f_{\mathrm c} = \frac{1}{2\pi RC}.
  \label{eq:rc}
\end{equation}

Při uvádění výsledků nestačí samotný výpočet; stručně vysvětlete,
co znamenají a~zda odpovídají stanoveným požadavkům.
